\documentclass[11pt]{article}
\usepackage[margin=1in]{geometry}
\usepackage{amsmath,amssymb,amsthm}
\usepackage{hyperref}
\usepackage[numbers,sort&compress]{natbib}

\newtheorem{theorem}{Theorem}
\newtheorem{lemma}{Lemma}
\newtheorem{corollary}{Corollary}
\newtheorem{definition}{Definition}
\newtheorem{remark}{Remark}

\newcommand{\ste}[1]{\texttt{#1}}
\newcommand{\cech}{\v{C}ech}

\title{Obstruction Size versus Representation Blow-Up:\\
Rectangle Covers, Fooling Sets, and the Refutation of an Equality\\
\large A literature review and annotated bibliography, oriented toward
the Lean mechanization in \ste{Ste.RepresentationBounds}}
\author{Project \texttt{thejeb}}
\date{July 24, 2026}

\begin{document}
\maketitle

\begin{abstract}
The project's open-question queue carried the item ``obstruction size
$=$ representation blow-up'': the conjecture that the concrete \cech{}
obstruction number of \ste{Ste.CechObstruction} --- the count of
compatible-but-stuck families of local sections for the singleton cover
--- \emph{equals} the representational cost that a coupled constraint
forces on any variable-separable (rectangular) encoding. This note
settles the item as a \emph{negative} result with a precise positive
replacement. There are two readings of ``blow-up''. Under the
\emph{box-gap} reading the equality is true but tautological: the
obstruction is by definition $|\mathrm{box}(T)|-|T|$, the overshoot of
the smallest single rectangle containing $T$. Under the
\emph{complexity} reading --- the minimum number $\rho(T)$ of
rectangles whose union is exactly $T$, the direct analogue of the
nondeterministic cover number of communication complexity --- the
equality is \emph{false in both directions}, and what is true is a
sandwich of bounds: $\mathrm{fool}(T)\le\rho(T)\le|T|$, with
$\rho(T)=1$ exactly when the \cech{} obstruction vanishes, and
$\rho(\mathrm{allEqual}_{n,k})=k$ exactly while the obstruction is
$k^n-k$. All of these statements, including both strict-separation
witnesses refuting the equality, are machine-checked in Lean~4 in
\ste{Ste.RepresentationBounds}. This review surveys the communication
complexity and extension-complexity literature that the definitions are
drawn from, and states precisely which fragments are mechanized.
Following the project rule --- distrust results that are not machine
proven --- all claims about the mechanization are limited to statements
that compile in CI; everything else is attributed to the cited
literature.
\end{abstract}

\section{The question, and why it is two questions}

Let $V$ be a set of variables with value types $(A_v)_{v\in V}$ and let
$T\subseteq\prod_{v}A_v$ be a constraint (a hypothesis set in the
set-theoretic-estimation reading). \ste{Ste.CechObstruction} computes,
for the singleton cover of $V$, the set of \emph{compatible families}
of local sections --- which, because singleton contexts have empty
overlaps, is exactly the box
\[
\mathrm{box}(T)\;=\;\prod_{v\in V}\pi_v(T),
\]
the product of the per-variable projections --- and defines the
obstruction number
\[
\mathrm{obs}(T)\;=\;\bigl|\mathrm{box}(T)\bigr|-\bigl|T\bigr|
\quad\in\;\mathbb{N}\cup\{\infty\},
\]
the count of phantom assignments that margin-based reasoning must
entertain and the coupling then kills. The open item asked whether
$\mathrm{obs}(T)$ \emph{equals} the ``forced representation blow-up''
of $T$. The answer depends entirely on what a representation is.

\paragraph{Reading 1: blow-up of the single-rectangle representation.}
If a ``representation'' is one rectangle (one product set) and
``blow-up'' is how many spurious points the best such representation
adds, then the best rectangle containing $T$ is $\mathrm{box}(T)$
itself, and the blow-up is $|\mathrm{box}(T)|-|T| = \mathrm{obs}(T)$
\emph{by definition}. The Lean statement
(\ste{cechObstruction\_eq\_box\_sub\_encard}, with
\ste{compatibleFamilies\_eq\_boxProduct} exhibiting the box) is proved
by unfolding two definitions. The equality holds, but it is
bookkeeping, not a representation theorem; nothing in it deserves the
name ``blow-up equality''.

\paragraph{Reading 2: blow-up as the size of the smallest exact
rectangular representation.}
The substantive reading, standard since \citet{aho1983notions} and
systematized by \citet{kushilevitz1997communication}, measures a
constraint by its \emph{rectangle-cover complexity}
\[
\rho(T)\;=\;\min\Bigl\{m \;:\; T=R_1\cup\dots\cup R_m,\ \ R_i
\text{ rectangles}\Bigr\},
\]
the minimum number of product sets whose union is exactly $T$ (each
$R_i\subseteq T$ automatically). In two-party communication complexity
$\log_2\rho$ is the nondeterministic complexity of the predicate $T$ up
to one bit; in the extension-complexity literature
\citep{yannakakis1991expressing,fiorini2015exponential} $\rho$ is the
combinatorial lower bound (the rectangle covering bound) on
nonnegative rank and hence on the size of linear-programming
reformulations. Under this reading the conjectured equality
$\mathrm{obs}(T)=\rho(T)$ (or any uniform inequality between the two)
is \emph{false}, and the correct picture is the sandwich of
Section~\ref{sec:sandwich}.

\section{The sandwich, as mechanized}\label{sec:sandwich}

All results in this section are machine-checked, \ste{sorry}-free, in
\ste{lean/Ste/RepresentationBounds.lean}; names in typewriter font are
the Lean identifiers. $\rho$ is \ste{rectCoverNumber}, valued in
$\mathbb{N}\cup\{\infty\}$ ($\infty$ if no finite cover exists);
$\mathrm{fool}$ is \ste{foolingNumber}, the supremum of cardinalities
of fooling sets.

\begin{definition}[fooling set; \ste{IsFoolingSet}]
$F\subseteq T$ is a \emph{fooling set} for $T$ if no rectangle
$R\subseteq T$ contains two distinct points of $F$.
\end{definition}

This is the abstract form of the classical definition: the usual
condition --- for $f\ne g$ in $F$ some coordinate mixture of $f$ and
$g$ leaves $T$ --- implies it, because rectangles are closed under
coordinatewise mixing (\ste{IsRectangle.mix\_mem}), which is the
mechanized engine of the argument.

\begin{theorem}[lower bound;
\ste{IsFoolingSet.encard\_le\_of\_hasRectCover},
\ste{foolingNumber\_le\_rectCoverNumber}]
If $F$ is a fooling set for $T$, then $|F|\le\rho(T)$; hence
$\mathrm{fool}(T)\le\rho(T)$. Each rectangle of an exact cover is a
sub-rectangle of $T$ and therefore contains at most one fooling point;
the mechanized proof builds the induced injection $F\to\{1,\dots,m\}$.
\end{theorem}

\begin{theorem}[upper bound; \ste{rectCoverNumber\_le\_encard}]
$\rho(T)\le|T|$: singletons are rectangles
(\ste{isRectangle\_singleton}), so $T$ covers itself pointwise. (For
infinite $T$ the right side is $\infty$ and the bound is vacuous.)
\end{theorem}

\begin{theorem}[bottom of the scale;
\ste{rectCoverNumber\_eq\_one\_iff},
\ste{rectCoverNumber\_eq\_one\_iff\_cechVanishes},
\ste{one\_lt\_rectCoverNumber\_iff}]
$\rho(T)=1$ iff $T$ is a nonempty rectangle; for nonempty $T$,
$\rho(T)=1$ iff the \cech{} obstruction vanishes, and $\rho(T)\ge 2$
iff it does not. (Via \ste{cechVanishes\_iff\_rectangular} of
\ste{Ste.CouplingLowerBound}.) At the bottom of the scale, obstruction
and cover complexity carry the same information --- this, and only
this, survives of the conjectured equality.
\end{theorem}

\begin{theorem}[the coupling, exactly;
\ste{isFoolingSet\_allEqual}, \ste{rectCoverNumber\_allEqual},
\ste{foolingNumber\_allEqual}]
For the $n$-fold coupling $\mathrm{allEqual}_{n,\alpha}$ (all
coordinates equal over an alphabet of size $k$) with $n\ge 2$: the $k$
constant vectors form a fooling set, and the sandwich pinches:
$\mathrm{fool}=\rho=k$. Meanwhile
$\mathrm{obs}(\mathrm{allEqual}_{n,\alpha})=k^n-k$
(\ste{cechObstruction\_allEqual}, proved previously).
\end{theorem}

\begin{theorem}[refutation of the equality, both directions;
\ste{rectCoverNumber\_lt\_cechObstruction\_allEqual},
\ste{exists\_rectCoverNumber\_lt\_cechObstruction},
\ste{exists\_cechObstruction\_lt\_rectCoverNumber},
\ste{not\_forall\_cechObstruction\_le\_rectCoverNumber},
\ste{not\_forall\_rectCoverNumber\_le\_cechObstruction}]
No uniform inequality holds between $\mathrm{obs}$ and $\rho$:
\begin{itemize}
\item every nonempty rectangle has $\mathrm{obs}=0<1=\rho$ (witness:
the full square on one Boolean variable), so
$\rho\le\mathrm{obs}$ fails;
\item for $n\ge3$, $k\ge2$,
$\rho(\mathrm{allEqual}_{n,\alpha})=k<k^n-k=\mathrm{obs}$ (witness:
$\mathrm{allEqual}_{3,\mathrm{Bool}}$ with $2<6$), so
$\mathrm{obs}\le\rho$ fails.
\end{itemize}
\end{theorem}

\begin{remark}[the origin of the conjecture;
\ste{diagonal\_eq\_allEqual}, \ste{rectCoverNumber\_diagonal},
\ste{diagonal\_cechObstruction\_eq\_rectCoverNumber}]
At the minimal nonvanishing instance --- the two-variable Boolean
diagonal, which is $\mathrm{allEqual}_{2,\mathrm{Bool}}$ --- the two
invariants \emph{coincide}: $\mathrm{obs}=2^2-2=2=\rho$. The equality
conjecture is the smallest example mistaken for a law: it fails at
$n=3$ ($6$ vs.\ $2$) and in the opposite direction for every nonempty
rectangle ($0$ vs.\ $1$).
\end{remark}

The conceptual moral matches the communication-complexity literature:
$\mathrm{obs}$ is a \emph{volume} invariant (how much the box
overshoots), $\rho$ is a \emph{covering} invariant (how many boxes are
needed); volume bounds covering only through logarithms, never linearly.
The classical relation --- $\log_2\rho$ equals nondeterministic
communication complexity up to an additive bit, and
$\mathrm{fool}(T)$ can be exponentially smaller than $\rho(T)$
\citep{dietzfelbinger1996comparison} --- lives one abstraction level
above what is mechanized here and remains outlook.

\section{Annotated bibliography}

\begin{description}

\item[\citet{kushilevitz1997communication}]
The standard monograph on two-party communication complexity; Chapters
1--2 define combinatorial rectangles, monochromatic covers, the cover
numbers $C^0,C^1$, and the fooling-set and rank lower-bound methods.
Our $\rho(T)$ is precisely their $C^1$ read on the constraint $T$ as a
$|V|$-dimensional 0/1 tensor (the book treats $|V|=2$; the rectangle
and fooling formalism, and our Lean proofs, are arity-uniform). The
book's Proposition~2.2 --- every cover rectangle contains at most one
fooling point --- is exactly the statement mechanized as
\ste{IsFoolingSet.encard\_le\_of\_hasRectCover}; the book's
$\log$-scale accounting ($N^1(f)=\log_2 C^1(f)$) is the part we
deliberately do not mechanize, since the STE obstruction is counted on
the linear scale.

\item[\citet{aho1983notions}]
The paper that introduced rectangle covers and the fooling-set bound in
the VLSI/communication setting, and posed the question (later resolved
negatively) of whether cover number and deterministic complexity are
polynomially related. Historically this is where ``representation size
of a relation $=$ number of monochromatic rectangles'' enters the
literature, which is why we take it as the anchor for the complexity
reading of ``blow-up''. The fooling-set method appears here in the
concrete mixture form whose abstract closure property is our
\ste{IsRectangle.mix\_mem}.

\item[\citet{yannakakis1991expressing}]
Connects rectangle covers to representation size in the strongest
classical sense: the size of the smallest linear program expressing a
combinatorial polytope (its extension complexity) is governed by the
nonnegative rank of a slack matrix, of which the rectangle covering
number is the support lower bound; symmetric LPs for matching and TSP
are shown exponential. For the STE reading, this is the license to call
$\rho$ a ``representation blow-up'': a constraint with $\rho(T)=m$ is
exactly a constraint expressible as a disjunction of $m$
variable-separable pieces and no fewer. Our
$\rho(\mathrm{allEqual})=k$ theorem is a (very small) instance of the
slack-support paradigm: the fooling set plays the role of a
combinatorially independent family in the slack matrix.

\item[\citet{dietzfelbinger1996comparison}]
Quantifies the looseness of the fooling-set method against the rank
method: fooling sets can be quadratically smaller than deterministic
complexity bounds, and (via a counting argument) almost all functions
have maximal complexity but only polynomial-size fooling sets. Relevant
here as a caution: our sandwich
$\mathrm{fool}(T)\le\rho(T)\le|T|$ pinches to equality on
$\mathrm{allEqual}$, but the left inequality is in general far from
tight, so \ste{foolingNumber} should not be mistaken for a
characterization of $\rho$. Any future STE conjecture of the form
``fooling number $=$ cover number for coupling-like constraints'' must
be tested against the separations in this paper.

\item[\citet{fiorini2015exponential}]
The modern culmination of Yannakakis's program: unconditional
exponential lower bounds on the extension complexity of the cut, stable
set, and TSP polytopes, obtained by lower-bounding the rectangle
covering/nonnegative-rank of explicit slack matrices (building on the
correlation polytope and quantum communication arguments). For the STE
project this fixes the right asymptotic vocabulary --- obstructions to
small rectangular representations are proved by exhibiting structured
sub-matrices, not by counting phantom assignments --- reinforcing the
present verdict that $\mathrm{obs}$, a phantom count, cannot equal a
covering complexity in general.

\item[\citet{abramsky2011sheaf}]
The sheaf-theoretic frame the STE \cech{} obstruction is modeled on:
empirical models over a cover of contexts, with non-contextuality
exactly the existence of a global section gluing the local ones. Our
$\mathrm{obs}(T)$ is the crudest numerical shadow of their obstruction
--- the count of compatible-but-stuck families for the singleton cover
--- and the mechanized equivalence ``$\mathrm{obs}$ vanishes iff $T$ is
a rectangle iff $\rho(T)=1$'' is the STE analogue of their
``non-contextual iff product structure'' at the coarsest cover. The
paper's cohomological refinements (with \ste{Ste.CechCover},
\ste{Ste.CechComplex}, \ste{Ste.TwistedCech}) are where a
\emph{graded} obstruction-vs-cover-number comparison would have to
live; nothing at that level is claimed here.

\item[\citet{dechter2003constraint}]
The constraint-processing side of the same divide: bounded induced
width buys a linear-size faithful junction-tree representation
(mechanized earlier in \ste{Ste.JunctionTree}), while
$\mathrm{allEqual}$, of maximal width, is exactly the constraint whose
margins carry no information. The present note sharpens the hard side
of that dichotomy: the coupling's true rectangular representation cost
is $k$ (small!), and what explodes is not $\rho$ but the box gap
$k^n-k$ --- the price of \emph{margin-only} reasoning, not of
\emph{rectangular} reasoning. This distinction was invisible while the
open item conflated the two.

\end{description}

\section{Verdict and open questions}

\textbf{Verdict.} The open item ``obstruction size $=$ representation
blow-up'' is \emph{closed, negatively}, with the following
machine-checked replacement. (i)~As a box-gap statement the equality is
a definition, not a theorem
(\ste{cechObstruction\_eq\_box\_sub\_encard}). (ii)~As a statement
about rectangle-cover complexity it is false in both directions:
$\mathrm{obs}=0<\rho=1$ on nonempty rectangles, and
$\rho=k<k^n-k=\mathrm{obs}$ on $\mathrm{allEqual}_{n,k}$ for $n\ge3$,
$k\ge2$
(\ste{not\_forall\_*}). (iii)~What is true and now mechanized is the
sandwich $\mathrm{fool}(T)\le\rho(T)\le|T|$, the boundary agreement
$\rho(T)=1\Leftrightarrow\mathrm{obs}(T)$ vanishes, and the exact value
$\rho(\mathrm{allEqual}_{n,\alpha})=|\alpha|$ via the constant-vector
fooling set. The equality holds at exactly one place we know of ---
the $2\times2$ diagonal, $\mathrm{obs}=\rho=2$ --- and that coincidence
is the likely origin of the conjecture.

\textbf{Open questions for the queue.}
(1)~\emph{Logarithmic reconciliation:} is there a uniform bound
$\rho(T)\le f(\mathrm{obs}(T),|T|)$, e.g.\ polynomial in
$\log(\mathrm{obs}+2)$ and $\log(|T|+2)$, for finite constraints? The
$\mathrm{allEqual}$ family is consistent with
$\rho\le\mathrm{obs}+1$ for \emph{nonempty non-rectangular} $T$
--- is $\rho(T)\le\mathrm{obs}(T)+1$ always true there? (The fooling
bound gives no obstruction to it; a proof or counterexample is a
concrete next target.)
(2)~\emph{Partition versus cover:} the disjoint-cover (partition)
number $\rho^{\sqcup}$ satisfies
$\rho\le\rho^{\sqcup}$; does $\mathrm{allEqual}$ also have
$\rho^{\sqcup}=k$ (the singleton-of-constants partition gives
$\le k$, and the fooling bound gives $\ge k$ --- this should mechanize
with the same machinery).
(3)~\emph{Graded comparison:} for covers with nonempty overlaps
(\ste{Ste.CechCover}) is there a cover-relative $\rho$ for which the
vanishing boundary ($\rho=1$) again matches the cover-level sheaf
condition?

\bibliographystyle{plainnat}
\bibliography{../../refs}

\end{document}
